%% ma: L12-B06-0080
%% lop: 12
%% chuong: 2
%% bai: 6
%% dang: 12.6.21
%% loai: TN4
%% muc: TH
%% dap_an: "A"
%% nguon: "CĐ2.6.pdf giáo viên gửi 27/10/2026 (M07, Lê Minh Khôi – Nguyễn Anh Khôi, Chương II Chủ đề 6 Ứng dụng vào các bài toán thực tế), Câu 4"
%% kien_thuc: "Bài 6 (tích của vectơ với một số, hai vectơ cùng hướng)"
%% trang_thai: chinh-thuc
%% ghi_chu: "Giáo viên đã duyệt 27/10/2026. Bỏ hình minh họa máy bay và bốn lực (hình đời sống, không chứa số liệu)"
\begin{ex}
Khi chuyển động trong không gian, máy bay luôn chịu tác động của bốn lực chính: lực đẩy của động cơ, lực cản của không khí, trọng lực và lực nâng khí động học. Lực cản của không khí ngược hướng với lực đẩy của động cơ và có độ lớn tỉ lệ thuận với bình phương vận tốc máy bay. Một chiếc máy bay tăng vận tốc từ $900$ km/h lên $920$ km/h, trong quá trình tăng tốc máy bay giữ nguyên hướng bay. Lực cản của không khí khi máy bay đạt vận tốc $900$ km/h và $920$ km/h lần lượt được biểu diễn bởi hai vectơ $\vec F_1$ và $\vec F_2$. Biết $\vec F_1=k\cdot\vec F_2$ với $k$ là một số thực dương. Tính giá trị của $k$ (làm tròn kết quả đến hàng phần trăm).
\choice{\True $k\approx0{,}96$}{$k\approx1{,}04$}{$k\approx0{,}98$}{$k\approx0{,}90$}
\loigiai{Hai lực cùng hướng, $k=\dfrac{|\vec F_1|}{|\vec F_2|}=\left(\dfrac{900}{920}\right)^2\approx0{,}96$.}
\end{ex}
